\documentclass[11pt, a4paper]{article}

\usepackage[a4paper, margin=2cm]{geometry}
\usepackage[spanish]{babel}
\usepackage{lmodern}
\usepackage[scaled=0.95]{helvet}
\usepackage{courier}

\usepackage{amsmath}
\usepackage{booktabs}
\usepackage{xcolor}
\usepackage{graphicx}
\usepackage{listings}
\usepackage{tcolorbox}
\tcbuselibrary{skins, breakable, listings}
\usepackage{hyperref}

% ==========================================
% PALETA DE COLORES DE ALTO CONTRASTE
% ==========================================
\definecolor{brandNavy}{RGB}{15, 23, 42}       % Deep Navy Dark
\definecolor{brandBlue}{RGB}{37, 99, 235}      % Vibrant Blue Accent
\definecolor{brandCyan}{RGB}{14, 165, 233}     % Cyan Light Accent
\definecolor{codeBg}{RGB}{248, 250, 252}        % Canvas Soft Light Gray
\definecolor{codeFrame}{RGB}{203, 213, 225}     % Border Slate Gray
\definecolor{codeComment}{RGB}{100, 116, 139}   % Muted Comment Text
\definecolor{codeKeyword}{RGB}{3, 105, 161}     % Deep Teal-Blue Keyword
\definecolor{codeString}{RGB}{180, 83, 9}       % Warm Amber String
\definecolor{termBg}{RGB}{15, 23, 42}          % Dark Terminal Canvas
\definecolor{termBorder}{RGB}{30, 41, 59}      % Terminal Border
\definecolor{termFg}{RGB}{241, 245, 249}        % Terminal Text Light
\definecolor{boxTitleBg}{RGB}{30, 58, 138}     % Box Header Dark Blue

% ==========================================
% CONFIGURACIÓN DE LISTINGS Y TCOLORBOX
% ==========================================
\lstdefinestyle{codigo}{
  language=C++,
  basicstyle=\small\ttfamily\color{brandNavy},
  keywordstyle=\bfseries\color{codeKeyword},
  stringstyle=\color{codeString},
  commentstyle=\itshape\color{codeComment},
  numbers=left,
  numberstyle=\tiny\bfseries\color{codeComment},
  stepnumber=1,
  numbersep=10pt,
  showstringspaces=false,
  breaklines=true,
  tabsize=2
}

\newtcblisting{codebox}[1][]{
  toptitle=2mm,
  bottomtitle=2mm,
  colback=codeBg,
  colframe=brandBlue,
  coltitle=white,
  fonttitle=\bfseries\small\sffamily,
  enhanced,
  attach boxed title to top left={xshift=5mm, yshift=-3mm},
  boxed title style={colback=brandBlue, arc=2mm},
  arc=3mm,
  boxrule=1.2pt,
  breakable,
  listing only,
  listing engine=listings,
  listing options={style=codigo},
  #1
}

\newtcblisting{terminalout}[1][]{
  toptitle=2mm,
  bottomtitle=2mm,
  colback=termBg,
  colframe=termBorder,
  colupper=termFg,
  fontupper=\small\ttfamily,
  fonttitle=\bfseries\small\sffamily,
  title={> COMPILACIÓN Y EJECUCIÓN EN CONSOLA},
  coltitle=white,
  enhanced,
  attach boxed title to top left={xshift=5mm, yshift=-3mm},
  boxed title style={colback=boxTitleBg, arc=2mm},
  arc=3mm,
  boxrule=1.2pt,
  breakable,
  listing only,
  listing options={
    basicstyle=\small\ttfamily\color{termFg},
    breaklines=true,
    showstringspaces=false
  },
  #1
}

\hypersetup{
    colorlinks=true,
    linkcolor=brandBlue,
    urlcolor=brandBlue,
    pdftitle={Ejercicios de Pilas (Stack) en C++ - FINESI},
    pdfauthor={Fred Torres Cruz}
}

\begin{document}

% ==========================================
% PORTADA CON LOGOS
% ==========================================
\begin{titlepage}
    \centering
    
    % Logos alineados horizontalmente con comprobación de existencia
    \begin{minipage}{0.42\textwidth}
        \begin{center}
            \IfFileExists{logo-una-puno.png}{%
                \includegraphics[height=3.2cm, keepaspectratio]{logo-una-puno.png}%
            }{%
                \begin{tcolorbox}[colback=brandNavy, colframe=brandBlue, width=3.2cm, height=3.2cm, arc=2mm, halign=center, valign=center, boxrule=1pt]
                    {\small \bfseries \color{white} LOGO\\UNA-PUNO}
                \end{tcolorbox}%
            }
        \end{center}
    \end{minipage}
    \hfill
    \begin{minipage}{0.42\textwidth}
        \begin{center}
            \IfFileExists{logofinesi.png}{%
                \includegraphics[height=3.2cm, keepaspectratio]{logofinesi.png}%
            }{%
                \begin{tcolorbox}[colback=brandNavy, colframe=brandBlue, width=3.2cm, height=3.2cm, arc=2mm, halign=center, valign=center, boxrule=1pt]
                    {\small \bfseries \color{white} LOGO\\FINESI}
                \end{tcolorbox}%
            }
        \end{center}
    \end{minipage}

    \vspace{1.2cm}

    {\Large \bfseries UNIVERSIDAD NACIONAL DEL ALTIPLANO \par}
    {\large FACULTAD DE INGENIERÍA ESTADÍSTICA E INFORMÁTICA \par}

    \vspace{1.2cm}

    \begin{tcolorbox}[colback=brandNavy, colframe=brandBlue, arc=4mm, width=\textwidth, halign=center, top=8mm, bottom=8mm]
        {\huge \bfseries \color{white} EJERCICIOS DE PILAS (STACK) \par}
        \vspace{0.3cm}
        {\Large \color{brandCyan} Práctica Introductoria -- Principio LIFO y Operaciones Básicas \par}
    \end{tcolorbox}
    
    \vspace{1.0cm}
    
    {\huge \bfseries \color{brandNavy} ESTRUCTURA DE DATOS \par}
    \vspace{0.4cm}
    \rule{0.8\textwidth}{1.5pt}
    \vspace{0.8cm}
    
    {\Large \textbf{Docente:} \color{brandBlue} Fred Torres Cruz \par}
    {\Large \textbf{ Estudiante:} \color{brandBlue} Blays Kevin Mamani Quispe \par} {\Large \textbf{ Link de Repositorio GitHub:} \color{brandBlue} https://github.com/blayskevinmamaniquispe-lab/Estructura-de-Datos/tree/main/Ejerccicios%20-%20Estructuras%20b%C3%A1sicas. \par} \vfill
    \vfill
    
    \begin{tcolorbox}[colback=codeBg, colframe=brandBlue, arc=3mm, width=0.95\textwidth, halign=left, boxrule=1pt]
        \textbf{\large Indicaciones Generales:}
        \begin{itemize}\setlength{\itemsep}{3pt}
            \item Resuelva los ejercicios utilizando la estructura \texttt{std::stack} de C++.
            \item Identifique en todo momento el elemento ubicado en el \textbf{TOPE}.
            \item Observe y compruebe el comportamiento \textbf{LIFO} (\textit{Last In, First Out}).
        \end{itemize}
    \end{tcolorbox}
    
    \vfill
    
    {\large \bfseries \color{codeComment} Puno, Perú \par}
    {\large \color{codeComment} 2026 \par}
\end{titlepage}

\newpage
\tableofcontents
\newpage

% ==========================================
% CONCEPTO CLAVE Y PLANTILLA BASE
% ==========================================
\section*{Introducción: El Principio LIFO y Plantilla Base}
\addcontentsline{toc}{section}{Introducción: El Principio LIFO y Plantilla Base}

En una estructura de datos tipo \textbf{Pila (Stack)}, los elementos se agregan y remueven siguiendo el principio \textbf{LIFO} (\textit{Last In, First Out} -- Último en entrar, primero en salir).

\begin{center}
    \begin{tcolorbox}[colback=brandNavy, colframe=brandBlue, width=0.7\textwidth, halign=center]
        {\Large \color{white} \textbf{TOPE} $\rightarrow$ \textbf{30} \quad 20 \quad 10} \par
        \vspace{2mm}
        {\small \color{brandCyan} Idea clave: si se insertaron 10, 20 y 30, el primero en salir será 30.}
    \end{tcolorbox}
\end{center}

\subsection*{Plantilla Base en C++}

El siguiente código constituye la estructura mínima con \texttt{std::stack}:

\begin{codebox}[title={PLANTILLA BASE EN C++}]
#include <iostream>
#include <stack>
using namespace std;

int main() {
    stack<int> pila;

    pila.push(10);
    pila.push(20);
    pila.push(30);

    while (!pila.empty()) {
        cout << pila.top() << " ";
        pila.pop();
    }

    cout << endl;
    return 0;
}
\end{codebox}

\begin{terminalout}
$ g++ plantilla.cpp -o plantilla
$ ./plantilla
30 20 10
\end{terminalout}

\newpage

% ==========================================
% EJERCICIO 1
% ==========================================
\section*{Ejercicio 1. Apilar números enteros}
\addcontentsline{toc}{section}{Ejercicio 1. Apilar números enteros}

\textbf{Enunciado:} Crear un programa que permita ingresar 5 números enteros y almacenarlos en una pila. Luego, retirar y mostrar los elementos uno por uno hasta dejar la pila vacía. El orden de salida debe permitir comprobar el funcionamiento LIFO.

\textbf{Condición:} Use las operaciones \texttt{push()}, \texttt{top()}, \texttt{pop()} y \texttt{empty()}.

\begin{codebox}[title={CÓDIGO FUENTE: Ejercicio 1}]
#include <iostream>
#include <stack>

using namespace std;

int main() {
    stack<int> pila;
    int numero;

    cout << "--- INGRESO DE 5 NÚMEROS ENTEROS ---\n";
    for (int i = 1; i <= 5; i++) {
        cout << "Ingrese el numero " << i << ": ";
        cin >> numero;
        pila.push(numero);
    }

    cout << "\n--- DESAPILANDO ELEMENTOS (ORDEN LIFO) ---\n";
    while (!pila.empty()) {
        cout << "Elemento en el TOPE: " << pila.top() << endl;
        pila.pop();
    }

    return 0;
}
\end{codebox}

\begin{terminalout}
$ g++ ejercicio1.cpp -o ejercicio1
$ ./ejercicio1
--- INGRESO DE 5 NÚMEROS ENTEROS ---
Ingrese el numero 1: 4
Ingrese el numero 2: 8
Ingrese el numero 3: 12
Ingrese el numero 4: 16
Ingrese el numero 5: 20

--- DESAPILANDO ELEMENTOS (ORDEN LIFO) ---
Elemento en el TOPE: 20
Elemento en el TOPE: 16
Elemento en el TOPE: 12
Elemento en el TOPE: 8
Elemento en el TOPE: 4
\end{terminalout}

\newpage

% ==========================================
% EJERCICIO 2
% ==========================================
\section*{Ejercicio 2. Invertir una palabra}
\addcontentsline{toc}{section}{Ejercicio 2. Invertir una palabra}

\textbf{Enunciado:} Leer una palabra ingresada por el usuario. Insertar cada carácter en una pila y, posteriormente, desapilar los caracteres para construir y mostrar la palabra en orden inverso.

\textbf{Condición:} Analice por qué una pila permite invertir naturalmente el orden de los elementos.

\begin{codebox}[title={CÓDIGO FUENTE: Ejercicio 2}]
#include <iostream>
#include <stack>
#include <string>

using namespace std;

int main() {
    stack<char> pilaCaracteres;
    string palabra;

    cout << "Ingrese una palabra: ";
    cin >> palabra;

    // Apilar cada carácter de la palabra
    for (char c : palabra) {
        pilaCaracteres.push(c);
    }

    string palabraInvertida = "";

    // Desapilar caracteres para revertir el orden
    while (!pilaCaracteres.empty()) {
        palabraInvertida += pilaCaracteres.top();
        pilaCaracteres.pop();
    }

    cout << "\nPalabra original: " << palabra << endl;
    cout << "Palabra invertida: " << palabraInvertida << endl;

    return 0;
}
\end{codebox}

\begin{terminalout}
$ g++ ejercicio2.cpp -o ejercicio2
$ ./ejercicio2
Ingrese una palabra: DATOS

Palabra original: DATOS
Palabra invertida: SOTAD
\end{terminalout}

\newpage

% ==========================================
% EJERCICIO 3
% ==========================================
\section*{Ejercicio 3. Verificar paréntesis balanceados}
\addcontentsline{toc}{section}{Ejercicio 3. Verificar paréntesis balanceados}

\textbf{Enunciado:} Crear un programa que determine si una expresión tiene sus paréntesis correctamente balanceados. Cada vez que aparezca un paréntesis de apertura \texttt{(}, se almacena en la pila. Cuando aparezca \texttt{)}, se comprueba que exista una apertura pendiente.

\textbf{Condición:} La expresión será correcta únicamente si nunca se intenta cerrar un paréntesis inexistente y la pila queda vacía al finalizar.

\begin{codebox}[title={CÓDIGO FUENTE: Ejercicio 3}]
#include <iostream>
#include <stack>
#include <string>

using namespace std;

bool estaBalanceado(const string& expresion) {
    stack<char> pila;

    for (char c : expresion) {
        if (c == '(') {
            pila.push(c);
        } else if (c == ')') {
            if (pila.empty()) {
                return false; // Intento de cerrar paréntesis sin apertura
            }
            pila.pop();
        }
    }

    return pila.empty(); // Es correcto si no quedan aperturas pendientes
}

int main() {
    string expresion;

    cout << "Ingrese la expresion matematica: ";
    cin >> expresion;

    if (estaBalanceado(expresion)) {
        cout << "Resultado: La expresion esta CORRECTAMENTE balanceada." << endl;
    } else {
        cout << "Resultado: La expresion esta INCORRECTAMENTE balanceada." << endl;
    }

    return 0;
}
\end{codebox}

\begin{terminalout}
$ g++ ejercicio3.cpp -o ejercicio3
$ ./ejercicio3
Ingrese la expresion matematica: (a+b)*(c-d)
Resultado: La expresion esta CORRECTAMENTE balanceada.

$ ./ejercicio3
Ingrese la expresion matematica: ((a+b)
Resultado: La expresion esta INCORRECTAMENTE balanceada.
\end{terminalout}

\newpage

% ==========================================
% EJERCICIO 4
% ==========================================
\section*{Ejercicio 4. Historial de acciones -- Deshacer (Undo)}
\addcontentsline{toc}{section}{Ejercicio 4. Historial de acciones -- Deshacer (Undo)}

\textbf{Enunciado:} Simular la función Deshacer de un editor. Cada acción realizada por el usuario se almacena en una pila. Cuando se solicite deshacer, el programa debe mostrar y eliminar la última acción registrada.

\textbf{Condición:} Incluya una opción para mostrar cuál es la acción actualmente ubicada en el tope de la pila.

\begin{codebox}[title={CÓDIGO FUENTE: Ejercicio 4}]
#include <iostream>
#include <stack>
#include <string>

using namespace std;

int main() {
    stack<string> historial;
    int opcion;
    string accion;

    do {
        cout << "\n--- EDITOR DE TEXTO (DESHACER) ---\n";
        cout << "1. Realizar nueva accion\n";
        cout << "2. Ver accion en el TOPE\n";
        cout << "3. Deshacer (Undo)\n";
        cout << "4. Salir\n";
        cout << "Seleccione una opcion: ";
        cin >> opcion;
        cin.ignore();

        switch (opcion) {
            case 1:
                cout << "Ingrese la accion realizada: ";
                getline(cin, accion);
                historial.push(accion);
                cout << "Accion registrada con exito.\n";
                break;
            case 2:
                if (!historial.empty()) {
                    cout << "TOPE actual: \"" << historial.top() << "\"\n";
                } else {
                    cout << "El historial esta vacio.\n";
                }
                break;
            case 3:
                if (!historial.empty()) {
                    cout << "Deshaciendo accion: \"" << historial.top() << "\"\n";
                    historial.pop();
                } else {
                    cout << "No hay acciones para deshacer.\n";
                }
                break;
            case 4:
                cout << "Saliendo del simulador...\n";
                break;
            default:
                cout << "Opcion no valida.\n";
        }
    } while (opcion != 4);

    return 0;
}
\end{codebox}

\begin{terminalout}
$ g++ ejercicio4.cpp -o ejercicio4
$ ./ejercicio4
--- EDITOR DE TEXTO (DESHACER) ---
1. Realizar nueva accion
2. Ver accion en el TOPE
3. Deshacer (Undo)
4. Salir
Seleccione una opcion: 1
Ingrese la accion realizada: escribir texto

Seleccione una opcion: 1
Ingrese la accion realizada: eliminar palabra

Seleccione una opcion: 2
TOPE actual: "eliminar palabra"

Seleccione una opcion: 3
Deshaciendo accion: "eliminar palabra"
\end{terminalout}

\newpage

% ==========================================
% EJERCICIO 5
% ==========================================
\section*{Ejercicio 5. Conversión de decimal a binario}
\addcontentsline{toc}{section}{Ejercicio 5. Conversión de decimal a binario}

\textbf{Enunciado:} Leer un número entero decimal positivo y convertirlo a binario dividiendo repetidamente entre 2 y almacenando el residuo en la pila. Finalmente, desapilar los residuos para obtener el número en el orden correcto.

\textbf{Condición:} No utilice funciones directas de conversión a binario; la solución debe evidenciar el uso de la pila.

\begin{codebox}[title={CÓDIGO FUENTE: Ejercicio 5}]
#include <iostream>
#include <stack>

using namespace std;

int main() {
    stack<int> residuos;
    int numeroDecimal;

    cout << "Ingrese un numero entero decimal positivo: ";
    cin >> numeroDecimal;

    if (numeroDecimal < 0) {
        cout << "Por favor ingrese un numero positivo." << endl;
        return 1;
    }

    if (numeroDecimal == 0) {
        cout << "Binario: 0" << endl;
        return 0;
    }

    int aux = numeroDecimal;
    while (aux > 0) {
        residuos.push(aux % 2);
        aux /= 2;
    }

    cout << "Numero decimal: " << numeroDecimal << endl;
    cout << "Representacion en binario: ";
    while (!residuos.empty()) {
        cout << residuos.top();
        residuos.pop();
    }
    cout << endl;

    return 0;
}
\end{codebox}

\begin{terminalout}
$ g++ ejercicio5.cpp -o ejercicio5
$ ./ejercicio5
Ingrese un numero entero decimal positivo: 13
Numero decimal: 13
Representacion en binario: 1101
\end{terminalout}

\newpage

% ==========================================
% PREGUNTA DE CIERRE
% ==========================================
\section*{Pregunta de Cierre y Análisis}
\addcontentsline{toc}{section}{Pregunta de Cierre y Análisis}

\begin{tcolorbox}[colback=codeBg, colframe=brandBlue, arc=3mm, width=\textwidth, boxrule=1.2pt]
    \textbf{\Large ¿Qué característica común tienen los cinco problemas que hace apropiado utilizar una pila?}
\end{tcolorbox}

\vspace{0.4cm}

\textbf{Respuesta y Justificación Técnica:}

La característica común fundamental es el **requerimiento de un procesamiento en orden inverso al de llegada**, donde el **último elemento procesado o ingresado es el primero en ser atendido o recuperado** (Principio \textbf{LIFO} -- \textit{Last In, First Out}).

\begin{enumerate}\setlength{\itemsep}{6pt}
 \item **Inversión de Secuencias (Ejercicios 1 y 2):** Para desapilar números o una palabra de forma invertida, el carácter/número ingresado en último lugar debe salir primero.
    \item **Anidamiento y Pareo de Simetría (Ejercicio 3):** Al validar paréntesis, el último paréntesis de apertura encontrado es siempre el primero que debe cerrarse con uno de cierre.
    \item **Deshacer Operaciones en Cascada (Ejercicio 4):** En un historial de acciones, la última acción realizada por el usuario es siempre la primera que se debe revertir.
    \item **Reconstrucción del Orden Correcto (Ejercicio 5):** Los residuos binarios se generan del bit menos significativo (LSB) al más significativo (MSB); desapilarlos permite imprimir el número en la lectura correcta de izquierda a derecha.
\end{enumerate}

\end{document}